% Bản nháp chương 6 theo dàn bài mới, viết với bộ mẫu ở style/mau/.
% ===========================================================================
\chapter{Ba cơ chế giải thích mọi chênh lệch}
\label{sec:mechanisms}

Chương~\ref{sec:compare} để lại câu hỏi vì sao các đường lại nằm ở đúng những chỗ
đó. Nhóm đã quét bảy lưới tham số để trả lời, nhưng bảy lưới ấy không cho bảy bài
học rời rạc: mọi chênh lệch đo được trong báo cáo đều quy về ba cơ chế. Chương
này trình bày theo ba cơ chế đó thay vì theo bảy thuật toán, và mỗi mục dùng đúng
lưới tham số nào chiếu sáng cơ chế của nó.

\section{Cơ chế thứ nhất: độ dài bước so với ngưỡng \texorpdfstring{$2/L$}{2/L}}
\label{sec:mech-step}

\resultfig{gd_fixed_iter}
  {Gradient descent với các độ dài bước cố định, theo số vòng lặp.}
  {fig:gd-fixed-iter}
\resultfig{gd_fixed_time}
  {Gradient descent với các độ dài bước cố định, theo thời gian chạy.}
  {fig:gd-fixed-time}

\begin{table}[tp]
  \centering
  \caption{Gradient descent với độ dài bước cố định. Ngưỡng ổn định là
    $2/L = \num{0.745}$.}
  \label{tab:gd-fixed}
  \begin{tabular}{lrrl}
    \toprule
    Độ dài bước & Vòng lặp đạt $10^{-6}$ & Thời gian (s) & Kết cục \\
    \midrule
    $t = 2.1/L$      & không đạt & --           & phân kỳ, $f$ lên $5{,}4 \cdot 10^{11}$ \\
    $t = 2/L$        & 20        & \num{0.465}  & dừng ở $4{,}9 \cdot 10^{-8}$ \\
    $t = 1.9/L$      & 20        & \num{0.478}  & hội tụ \\
    $t = 2/(L+\mu)$  & 20        & \num{0.729}  & hội tụ \\
    $t = 1/L$        & 40        & \num{1.181}  & hội tụ \\
    $t = 0.5/L$      & 80        & \num{2.402}  & hội tụ \\
    $t = 0.1/L$      & 380       & \num{10.13}  & hội tụ \\
    $t = 0.01/L$     & không đạt & --           & 1200 vòng mới xuống $3{,}7 \cdot 10^{-4}$ \\
    \bottomrule
  \end{tabular}
\end{table}

Với $t = 2.1/L$ hàm mục tiêu tăng tới $5{,}4 \cdot 10^{11}$, còn với $t = 1.9/L$
phương pháp hội tụ sau 20 vòng lặp, dù hai độ dài bước chỉ cách nhau 10\%
(bảng~\ref{tab:gd-fixed}). Ngưỡng $2/L$ vì thế là một ranh giới thật chứ không
phải một cận an toàn nới rộng. Sai số của thành phần ứng với trị riêng
$\lambda_i$ của Hessian nhân với $\abs{1 - t\lambda_i}$ mỗi vòng lặp, nên hệ số
của thành phần ứng với $\lambda_{\max} = L$ vượt quá 1 ngay khi $t$ vượt $2/L$.
Đúng tại ngưỡng, hệ số đó bằng 1 nên thành phần này đứng yên, và đó là lý do
$t = 2/L$ dừng ở $4{,}9 \cdot 10^{-8}$ thay vì đi tiếp tới giới hạn số học.

Trong vùng ổn định, giảm $t$ đi một nửa thì số vòng lặp tăng gấp đôi, thấy rõ ở
dãy 20, 40, 80 ứng với $t = 1.9/L$, $1/L$ và $0.5/L$ trên
hình~\ref{fig:gd-fixed-iter}. Quan hệ tỉ lệ nghịch ấy chỉ đọc được khi ngân sách
vòng lặp còn đủ: $t = 0.01/L$ chạy hết 1200 vòng vẫn dừng ở
$3{,}7 \cdot 10^{-4}$ nên nó không cho một con số để đặt vào dãy trên. Thứ tự các
đường giữ nguyên khi đổi trục hoành sang thời gian chạy
(hình~\ref{fig:gd-fixed-time}), vì mọi cấu hình ở đây gọi gradient đúng một lần
mỗi vòng lặp nên thời gian chỉ là số vòng lặp nhân với một hằng số.

\subsection{Khi độ dài bước do line search chọn}
\label{sec:mech-backtracking}

\resultfig{gd_backtracking_iter}
  {Ảnh hưởng của $\alpha$ và $\beta$ tới gradient descent với backtracking, theo
   số vòng lặp.}
  {fig:gd-bt-iter}
\resultfig{gd_backtracking_time}
  {Cùng các cấu hình ấy theo thời gian chạy. Hai trục tách nhau ở đây, khác với
   hình~\ref{fig:gd-fixed-time}, vì số lần đánh giá hàm mỗi vòng lặp không còn
   giống nhau giữa các cấu hình.}
  {fig:gd-bt-time}

Backtracking bỏ được nhu cầu biết $L$, và cái giá phải trả nằm ở số lần đánh giá
hàm. Trong chín cấu hình đã quét, số lần đánh giá hàm mỗi vòng lặp chỉ phụ thuộc
$\beta$: nó bằng \num{2.40} với $\beta = 0.5$, khoảng \num{3.45} với
$\beta = 0.8$ và khoảng \num{5.1} với $\beta = 0.9$, gần như không đổi khi
$\alpha$ chạy từ \num{0.1} tới \num{0.5}. Điều đó khớp với vai trò của hai tham
số ở \eqref{eq:armijo}: $\beta$ quyết định mỗi lần thử thu nhỏ bước bao nhiêu nên
nó định số lần thử, còn $\alpha$ chỉ dịch ngưỡng chấp nhận nên nó hiếm khi đổi
được số lần thử thêm một đơn vị.

Chú ý rằng một trong ba con số ấy là chi phí bắt buộc: mỗi vòng lặp phải đánh giá
$f$ ít nhất một lần để so với điều kiện Armijo, nên phần thực sự do line search
sinh ra là \num{1.40}, \num{2.45} và \num{4.1} lần. Đây cũng là lý do
hình~\ref{fig:gd-bt-time} không còn là hình~\ref{fig:gd-bt-iter} vẽ lại trên
thang khác, khác hẳn cặp hình của bước cố định.

\subsection{Momentum phải khớp với độ dài bước thực tế}
\label{sec:mech-momentum}

\resultfig{agd_iter}
  {Các cấu hình momentum của phương pháp tăng tốc, theo số vòng lặp.}
  {fig:agd-iter}
\resultfig{agd_time}
  {Các cấu hình momentum theo thời gian chạy.}
  {fig:agd-time}

Ghép backtracking với momentum là chỗ cơ chế thứ nhất lộ ra rõ nhất, vì hai thành
phần ấy phải nói về cùng một độ dài bước. Công thức
$\beta = (\sqrt{\kappa}-1)/(\sqrt{\kappa}+1)$ chỉ đúng khi $t = 1/L$, trong khi
line search chấp nhận bước lớn hơn nhiều mỗi khi gradient nằm theo hướng phẳng.
Bản cài đặt đầu tiên của nhóm ghép bước do line search chọn với momentum tính cho
$1/L$, và hàm mục tiêu bùng lên $1{,}47 \cdot 10^{4}$. Nhóm viết lại thành
%
\begin{equation}
  \label{eq:momentum-from-step}
  \beta = \frac{1 - \sqrt{t\mu}}{1 + \sqrt{t\mu}} ,
\end{equation}
%
tức tính theo bước thực tế, và giữ một phép kiểm tra trong bộ kiểm thử để lỗi ấy
không quay lại.

Sửa \eqref{eq:momentum-from-step} vẫn chưa đủ. Trong ba cấu hình backtracking đã
quét, cấu hình $\alpha = 0.5$ với $t_0 = 1$ đạt $10^{-6}$ sau 20 vòng lặp và cấu
hình $\alpha = 0.3$ với $t_0 = 1/L$ đạt sau 35 vòng, còn cấu hình $\alpha = 0.3$
với $t_0 = 1$ thì không đạt và dừng ở $5{,}5 \cdot 10^{-5}$
(hình~\ref{fig:agd-iter}). Với $\alpha = 1/2$, điều kiện Armijo trở thành đúng bổ
đề giảm
%
\begin{equation}
  \label{eq:descent-lemma}
  f\bigl(y - t \nabla f(y)\bigr) \le f(y) - \frac{t}{2} \norm{\nabla f(y)}^{2} ,
\end{equation}
%
tức chính bất đẳng thức mà chứng minh hội tụ của phương pháp tăng tốc dựa vào.
Nói cách khác, với $\alpha < 1/2$ thì line search chấp nhận những bước mà
\eqref{eq:descent-lemma} không bảo đảm, và phần chứng minh mất chỗ dựa. Khởi tạo
$t_0 = 1/L$ là cách xử lý còn lại, vì khi ấy line search chỉ có thể thu nhỏ từ một
giá trị vốn đã an toàn.

Khởi động lại là phần cải thiện lớn nhất trong nhóm này và nó không tốn gì. Với
công thức $\beta_k = (k-1)/(k+2)$, thêm khởi động lại đưa tổng thời gian từ
\SI{6.67}{\second} còn \SI{2.97}{\second} mà số vòng lặp đạt $10^{-6}$ giữ nguyên
ở 20, vì momentum tăng dần gây dao động ở phần đuôi và điều kiện
$\nabla f(y_k)\trans (w_{k+1} - w_k) > 0$ cắt đúng các dao động ấy
(hình~\ref{fig:agd-time}).

\section{Cơ chế thứ hai: nhiễu của gradient ngẫu nhiên}
\label{sec:mech-noise}

\resultfig{sgd_common_step_epoch}
  {SGD với cùng một độ dài bước cho mọi kích thước lô, theo số lượt duyệt dữ
   liệu.}
  {fig:sgd-common-epoch}
\resultfig{sgd_common_step_time}
  {Cùng thí nghiệm ấy theo thời gian chạy.}
  {fig:sgd-common-time}

Hằng số Lipschitz của mất mát một mẫu là $\norm{x_i}^{2} + \lambda$, đo được là
\num{3109.7} so với $L = \num{2.683}$, tức lớn hơn 1159 lần. Hệ quả là một độ dài
bước an toàn cho gradient đầy đủ không an toàn cho gradient một mẫu. Dùng chung
$\eta = 0.1/L = \num{0.0373}$ cho mọi kích thước lô thì $B = 1$ phân kỳ hẳn, còn
sai số cuối giảm đều theo kích thước lô, từ $2{,}5 \cdot 10^{-2}$ với $B = 16$
xuống $1{,}1 \cdot 10^{-4}$ với $B = 1024$ (hình~\ref{fig:sgd-common-epoch} và
hình~\ref{fig:sgd-common-time}). Mỗi kích thước lô vì thế phải nhận độ dài bước
riêng, lấy theo hằng số trơn của lô
%
\begin{equation}
  \label{eq:minibatch-smoothness}
  L_B = L + \frac{L_{\max} - L}{B} ,
\end{equation}
%
nội suy giữa $L_{\max}$ khi $B = 1$ và $L$ khi $B = n$.

\subsection{Bước hằng dừng ở một lân cận}
\label{sec:mech-schedule}

\resultfig{sgd_schedule_epoch}
  {Bốn quy tắc chọn độ dài bước của SGD với $B = 64$, theo số lượt duyệt dữ liệu.}
  {fig:sgd-schedule-epoch}
\resultfig{sgd_schedule_time}
  {Bốn quy tắc ấy theo thời gian chạy.}
  {fig:sgd-schedule-time}

Bốn quy tắc ở \eqref{eq:schedules} chạy cùng $B = 64$, cùng 30 lượt duyệt dữ liệu
và cùng chi phí tính toán trong khoảng \SI{2.59}{\second} tới
\SI{2.65}{\second}, nhưng dừng ở bốn mức sai số khác nhau hẳn: bước hằng dừng ở
$9{,}6 \cdot 10^{-4}$, quy tắc $\eta_0/\sqrt{k+1}$ ở $1{,}5 \cdot 10^{-4}$, quy
tắc $\eta_0/(1+\gamma k)$ ở $2{,}9 \cdot 10^{-5}$, còn giảm theo bậc thang ở
$5{,}8 \cdot 10^{-6}$, tức tốt hơn bước hằng 165 lần với cùng ngân sách
(hình~\ref{fig:sgd-schedule-epoch}).

Cơ chế nằm ở chỗ phương sai của $g_k$ trong \eqref{eq:sgd} không tắt theo $k$.
Gần $\wstar$ thì gradient đầy đủ tiến về không nhưng gradient của một lô thì
không, nên bước cập nhật cuối cùng chỉ còn là nhiễu nhân với $\eta$, và nghiệm
dao động trong một lân cận có bán kính tỉ lệ với $\eta$. Muốn sai số tiếp tục
giảm thì phải để $\eta_k$ tắt thay, và ba quy tắc giảm dần khác nhau ở chỗ chúng
tắt nhanh chậm ra sao. Hình~\ref{fig:sgd-schedule-time} không thêm thông tin so
với hình~\ref{fig:sgd-schedule-epoch}, vì bốn quy tắc chỉ khác nhau ở một phép
nhân vô hướng nên chi phí mỗi lượt duyệt như nhau.

\subsection{Kích thước lô và thời gian chạy}
\label{sec:mech-batch}

\resultfig{sgd_batch_epoch}
  {SGD với các kích thước lô khác nhau, mỗi lô dùng độ dài bước $1/L_B$ riêng,
   theo số lượt duyệt dữ liệu.}
  {fig:sgd-batch-epoch}
\resultfig{sgd_batch_time}
  {Cùng thí nghiệm ấy theo thời gian chạy, nơi lô nhỏ phải trả giá cho số lần cập
   nhật.}
  {fig:sgd-batch-time}

Khi mỗi kích thước lô nhận độ dài bước $1/L_B$ riêng theo
\eqref{eq:minibatch-smoothness}, năm cấu hình cho sai số cuối gần như nhau, trải
từ $8{,}2 \cdot 10^{-4}$ tới $1{,}2 \cdot 10^{-3}$
(hình~\ref{fig:sgd-batch-epoch}). Chênh lệch dồn hết sang thời gian chạy:
$B = 1$ mất \SI{34.5}{\second} cho 30 lượt duyệt còn $B = 256$ chỉ mất
\SI{1.86}{\second}, tức nhanh hơn 18 lần với cùng số lần chạm vào dữ liệu
(hình~\ref{fig:sgd-batch-time}). Lô nhỏ trả giá không phải vì tính nhiều hơn mà
vì nó chia cùng khối lượng tính toán thành nhiều lần gọi hơn, và chi phí cố định
của mỗi lần gọi không giảm theo kích thước lô.

\section{Cơ chế thứ ba: chi phí mỗi vòng lặp}
\label{sec:mech-cost}

\resultfig{newton_iter}
  {Các biến thể của phương pháp Newton, theo số vòng lặp.}
  {fig:newton-iter}
\resultfig{newton_time}
  {Cùng các biến thể ấy theo thời gian chạy. Đây là trục duy nhất phân biệt được
   chúng.}
  {fig:newton-time}

Hai cơ chế trên nói về việc mỗi vòng lặp cắt được bao nhiêu sai số. Cơ chế thứ ba
nói về việc vòng lặp ấy đắt bao nhiêu, và nó là lý do thứ hạng trên trục vòng lặp
khác thứ hạng trên trục thời gian.

Ba biến thể Newton đều hội tụ sau đúng một vòng lặp theo
\eqref{eq:newton-one-step} nên hình~\ref{fig:newton-iter} không phân biệt được
chúng; chỉ hình~\ref{fig:newton-time} phân biệt được. Phân rã lại Hessian mỗi
vòng tốn \SI{0.136}{\second} còn dùng lại phân rã cũ tốn \SI{0.133}{\second}, tức
gần như không khác, đúng như phải thế khi cả hai chỉ chạy một vòng. Bản damped
thêm backtracking mất \SI{0.168}{\second} vì nó đánh giá $f$ ba lần thay vì hai,
và line search ở đây không có việc gì để làm: bước đầy đủ $t = 1$ luôn thỏa điều
kiện Armijo trên một hàm bậc hai.

Newton-CG thay phân rã Cholesky bằng gradient liên hợp nên mỗi bước chỉ cần một
tích Hessian nhân vector với chi phí $\bigO(nd)$, không bao giờ phải lập ma trận
$d \times d$. Đổi lại nó mất \SI{1.28}{\second} tới \SI{1.32}{\second}, tức chậm
hơn Newton khoảng 10 lần ở $d = 280$. Cân bằng chỉ đảo khi $d$ đủ lớn để số hạng
$\bigO(d^3)$ vượt hẳn chi phí của vài chục bước gradient liên hợp, và
$d = 280$ còn xa mức đó.

\resultfig{lbfgs_iter}
  {L-BFGS với các kích thước bộ nhớ, theo số vòng lặp.}
  {fig:lbfgs-iter}
\resultfig{lbfgs_time}
  {L-BFGS với các kích thước bộ nhớ, theo thời gian chạy.}
  {fig:lbfgs-time}

L-BFGS nằm giữa hai nhóm và bốn kích thước bộ nhớ đã thử đều đạt $10^{-6}$ sau
đúng 8 vòng lặp (hình~\ref{fig:lbfgs-iter}), trong khi tổng thời gian trải từ
\SI{1.38}{\second} với $m = 10$ tới \SI{1.97}{\second} với $m = 3$
(hình~\ref{fig:lbfgs-time}). Bộ nhớ lớn hơn không mua thêm vòng lặp nào vì bài
toán là hàm bậc hai và tám cặp hiệu đã đủ dựng lại phần phổ mà thuật toán cần,
còn chênh lệch thời gian giữa bốn cấu hình nằm trong khoảng dao động giữa các lần
chạy nên không nên đọc thành thứ hạng.

\section{Ba cơ chế, đọc lại bảng tổng hợp}
\label{sec:mech-summary}

Ba cơ chế trên giải thích đủ bảng~\ref{tab:summary} mà không cần thêm gì. Ba
phương pháp bậc một bám nhau ở phần đầu vì chúng cùng bị ngưỡng $2/L$ chặn và
cùng cắt các hướng có độ cong lớn với tốc độ như nhau, rồi tách nhau ở phần đuôi
vì chỉ momentum mới đổi được hệ số co rút khi $\mu$ lên tiếng. Hai cấu hình SGD
lệch nhau 87,1 đơn vị tiền không phải vì kích thước lô mà vì một bản để $\eta_k$
tắt còn bản kia không, tức vì cơ chế thứ hai. Và Newton dẫn đầu bảng dù chỉ chạy
một vòng lặp là vì cơ chế thứ ba: vòng lặp ấy đắt gấp nhiều lần một vòng gradient
nhưng vẫn rẻ hơn hai chục vòng gradient cộng lại.

Cả ba cơ chế đều được đo ở đúng một giá trị $\lambda$, tức ở đúng một số điều
kiện. Chương~\ref{sec:theory} đối chiếu chúng với tốc độ mà lý thuyết dự đoán, và
chương~\ref{sec:lambda} kiểm xem chúng có giữ nguyên khi $\kappa$ đổi hay không.
